\IfFileExists{stacks-project.cls}{%
\documentclass{stacks-project}
}{%
\documentclass{amsart}
}

% For dealing with references we use the comment environment
\usepackage{verbatim}
\newenvironment{reference}{\comment}{\endcomment}
%\newenvironment{reference}{}{}
\newenvironment{slogan}{\comment}{\endcomment}
\newenvironment{history}{\comment}{\endcomment}

% For commutative diagrams we use Xy-pic
\usepackage[all]{xy}

% We use 2cell for 2-commutative diagrams.
\xyoption{2cell}
\UseAllTwocells

% We use multicol for the list of chapters between chapters
\usepackage{multicol}

% This is generally recommended for better output
\usepackage{lmodern}
\usepackage[T1]{fontenc}

% For cross-file-references
\usepackage{xr-hyper}

% Package for hypertext links:
\usepackage{hyperref}

% For any local file, say "hello.tex" you want to link to please
% use \externaldocument[hello-]{hello}
\externaldocument[introduction-]{introduction}
\externaldocument[conventions-]{conventions}
\externaldocument[sets-]{sets}
\externaldocument[categories-]{categories}
\externaldocument[topology-]{topology}
\externaldocument[sheaves-]{sheaves}
\externaldocument[sites-]{sites}
\externaldocument[stacks-]{stacks}
\externaldocument[fields-]{fields}
\externaldocument[algebra-]{algebra}
\externaldocument[brauer-]{brauer}
\externaldocument[homology-]{homology}
\externaldocument[derived-]{derived}
\externaldocument[simplicial-]{simplicial}
\externaldocument[more-algebra-]{more-algebra}
\externaldocument[smoothing-]{smoothing}
\externaldocument[modules-]{modules}
\externaldocument[sites-modules-]{sites-modules}
\externaldocument[injectives-]{injectives}
\externaldocument[cohomology-]{cohomology}
\externaldocument[sites-cohomology-]{sites-cohomology}
\externaldocument[dga-]{dga}
\externaldocument[dpa-]{dpa}
\externaldocument[sdga-]{sdga}
\externaldocument[hypercovering-]{hypercovering}
\externaldocument[schemes-]{schemes}
\externaldocument[constructions-]{constructions}
\externaldocument[properties-]{properties}
\externaldocument[morphisms-]{morphisms}
\externaldocument[coherent-]{coherent}
\externaldocument[divisors-]{divisors}
\externaldocument[limits-]{limits}
\externaldocument[varieties-]{varieties}
\externaldocument[topologies-]{topologies}
\externaldocument[descent-]{descent}
\externaldocument[perfect-]{perfect}
\externaldocument[more-morphisms-]{more-morphisms}
\externaldocument[flat-]{flat}
\externaldocument[groupoids-]{groupoids}
\externaldocument[more-groupoids-]{more-groupoids}
\externaldocument[etale-]{etale}
\externaldocument[chow-]{chow}
\externaldocument[intersection-]{intersection}
\externaldocument[pic-]{pic}
\externaldocument[weil-]{weil}
\externaldocument[adequate-]{adequate}
\externaldocument[dualizing-]{dualizing}
\externaldocument[duality-]{duality}
\externaldocument[discriminant-]{discriminant}
\externaldocument[derham-]{derham}
\externaldocument[local-cohomology-]{local-cohomology}
\externaldocument[algebraization-]{algebraization}
\externaldocument[curves-]{curves}
\externaldocument[resolve-]{resolve}
\externaldocument[models-]{models}
\externaldocument[functors-]{functors}
\externaldocument[equiv-]{equiv}
\externaldocument[pione-]{pione}
\externaldocument[etale-cohomology-]{etale-cohomology}
\externaldocument[proetale-]{proetale}
\externaldocument[relative-cycles-]{relative-cycles}
\externaldocument[more-etale-]{more-etale}
\externaldocument[trace-]{trace}
\externaldocument[crystalline-]{crystalline}
\externaldocument[spaces-]{spaces}
\externaldocument[spaces-properties-]{spaces-properties}
\externaldocument[spaces-morphisms-]{spaces-morphisms}
\externaldocument[decent-spaces-]{decent-spaces}
\externaldocument[spaces-cohomology-]{spaces-cohomology}
\externaldocument[spaces-limits-]{spaces-limits}
\externaldocument[spaces-divisors-]{spaces-divisors}
\externaldocument[spaces-over-fields-]{spaces-over-fields}
\externaldocument[spaces-topologies-]{spaces-topologies}
\externaldocument[spaces-descent-]{spaces-descent}
\externaldocument[spaces-perfect-]{spaces-perfect}
\externaldocument[spaces-more-morphisms-]{spaces-more-morphisms}
\externaldocument[spaces-flat-]{spaces-flat}
\externaldocument[spaces-groupoids-]{spaces-groupoids}
\externaldocument[spaces-more-groupoids-]{spaces-more-groupoids}
\externaldocument[bootstrap-]{bootstrap}
\externaldocument[spaces-pushouts-]{spaces-pushouts}
\externaldocument[spaces-chow-]{spaces-chow}
\externaldocument[groupoids-quotients-]{groupoids-quotients}
\externaldocument[spaces-more-cohomology-]{spaces-more-cohomology}
\externaldocument[spaces-simplicial-]{spaces-simplicial}
\externaldocument[spaces-duality-]{spaces-duality}
\externaldocument[formal-spaces-]{formal-spaces}
\externaldocument[restricted-]{restricted}
\externaldocument[spaces-resolve-]{spaces-resolve}
\externaldocument[formal-defos-]{formal-defos}
\externaldocument[defos-]{defos}
\externaldocument[cotangent-]{cotangent}
\externaldocument[examples-defos-]{examples-defos}
\externaldocument[algebraic-]{algebraic}
\externaldocument[examples-stacks-]{examples-stacks}
\externaldocument[stacks-sheaves-]{stacks-sheaves}
\externaldocument[criteria-]{criteria}
\externaldocument[artin-]{artin}
\externaldocument[quot-]{quot}
\externaldocument[stacks-properties-]{stacks-properties}
\externaldocument[stacks-morphisms-]{stacks-morphisms}
\externaldocument[stacks-limits-]{stacks-limits}
\externaldocument[stacks-cohomology-]{stacks-cohomology}
\externaldocument[stacks-perfect-]{stacks-perfect}
\externaldocument[stacks-introduction-]{stacks-introduction}
\externaldocument[stacks-more-morphisms-]{stacks-more-morphisms}
\externaldocument[stacks-geometry-]{stacks-geometry}
\externaldocument[moduli-]{moduli}
\externaldocument[moduli-curves-]{moduli-curves}
\externaldocument[examples-]{examples}
\externaldocument[exercises-]{exercises}
\externaldocument[guide-]{guide}
\externaldocument[desirables-]{desirables}
\externaldocument[coding-]{coding}
\externaldocument[obsolete-]{obsolete}
\externaldocument[fdl-]{fdl}
\externaldocument[index-]{index}

% Theorem environments.
%
\theoremstyle{plain}
\newtheorem{theorem}[subsection]{Theorem}
\newtheorem{proposition}[subsection]{Proposition}
\newtheorem{lemma}[subsection]{Lemma}

\theoremstyle{definition}
\newtheorem{definition}[subsection]{Definition}
\newtheorem{example}[subsection]{Example}
\newtheorem{exercise}[subsection]{Exercise}
\newtheorem{situation}[subsection]{Situation}

\theoremstyle{remark}
\newtheorem{remark}[subsection]{Remark}
\newtheorem{remarks}[subsection]{Remarks}

\numberwithin{equation}{subsection}

% Macros
%
\def\lim{\mathop{\mathrm{lim}}\nolimits}
\def\colim{\mathop{\mathrm{colim}}\nolimits}
\def\Spec{\mathop{\mathrm{Spec}}}
\def\Hom{\mathop{\mathrm{Hom}}\nolimits}
\def\Ext{\mathop{\mathrm{Ext}}\nolimits}
\def\SheafHom{\mathop{\mathcal{H}\!\mathit{om}}\nolimits}
\def\SheafExt{\mathop{\mathcal{E}\!\mathit{xt}}\nolimits}
\def\Sch{\mathit{Sch}}
\def\Mor{\mathop{\mathrm{Mor}}\nolimits}
\def\Ob{\mathop{\mathrm{Ob}}\nolimits}
\def\Sh{\mathop{\mathit{Sh}}\nolimits}
\def\NL{\mathop{N\!L}\nolimits}
\def\CH{\mathop{\mathrm{CH}}\nolimits}
\def\proetale{{pro\text{-}\acute{e}tale}}
\def\etale{{\acute{e}tale}}
\def\QCoh{\mathit{QCoh}}
\def\Ker{\mathop{\mathrm{Ker}}}
\def\Im{\mathop{\mathrm{Im}}}
\def\Coker{\mathop{\mathrm{Coker}}}
\def\Coim{\mathop{\mathrm{Coim}}}

% Boxtimes
%
\DeclareMathSymbol{\boxtimes}{\mathbin}{AMSa}{"02}

%
% Macros for moduli stacks/spaces
%
\def\QCohstack{\mathcal{QC}\!\mathit{oh}}
\def\Cohstack{\mathcal{C}\!\mathit{oh}}
\def\Spacesstack{\mathcal{S}\!\mathit{paces}}
\def\Quotfunctor{\mathrm{Quot}}
\def\Hilbfunctor{\mathrm{Hilb}}
\def\Curvesstack{\mathcal{C}\!\mathit{urves}}
\def\Polarizedstack{\mathcal{P}\!\mathit{olarized}}
\def\Complexesstack{\mathcal{C}\!\mathit{omplexes}}
% \Pic is the operator that assigns to X its picard group, usage \Pic(X)
% \Picardstack_{X/B} denotes the Picard stack of X over B
% \Picardfunctor_{X/B} denotes the Picard functor of X over B
\def\Pic{\mathop{\mathrm{Pic}}\nolimits}
\def\Picardstack{\mathcal{P}\!\mathit{ic}}
\def\Picardfunctor{\mathrm{Pic}}
\def\Deformationcategory{\mathcal{D}\!\mathit{ef}}
\usepackage{cleveref}
\crefname{originallemma}{original lemma}{original lemmas}
\Crefname{originallemma}{Original lemma}{Original lemmas}
\makeatletter
\newlabel{varieties-lemma-kunneth}{{0BED}{source}{}{sourcepointer-1}{}}
\newlabel{varieties-lemma-kunneth@cref}{{[originallemma][0][]0BED}{[1][0][]source}{}{}{}}
\makeatother
\crefname{originalsection}{original section}{original sections}
\Crefname{originalsection}{Original section}{Original sections}
\makeatletter
\newlabel{cohomology-section-cech-cohomology-of-complexes}{{01FP}{source}{}{sourcepointer-2}{}}
\newlabel{cohomology-section-cech-cohomology-of-complexes@cref}{{[originalsection][0][]01FP}{[1][0][]source}{}{}{}}
\makeatother
\crefname{originalequation}{original equation}{original equations}
\Crefname{originalequation}{Original equation}{Original equations}
\makeatletter
\newlabel{cohomology-equation-needs-signs}{{07MB}{source}{}{sourcepointer-3}{}}
\newlabel{cohomology-equation-needs-signs@cref}{{[originalequation][0][]07MB}{[1][0][]source}{}{}{}}
\makeatother
\crefname{originallemma}{original lemma}{original lemmas}
\Crefname{originallemma}{Original lemma}{Original lemmas}
\makeatletter
\newlabel{cohomology-lemma-diagrams-commute}{{0FKV}{source}{}{sourcepointer-4}{}}
\newlabel{cohomology-lemma-diagrams-commute@cref}{{[originallemma][0][]0FKV}{[1][0][]source}{}{}{}}
\makeatother
\crefname{originallemma}{original lemma}{original lemmas}
\Crefname{originallemma}{Original lemma}{Original lemmas}
\makeatletter
\newlabel{coherent-lemma-quasi-coherent-affine-cohomology-zero}{{01XB}{source}{}{sourcepointer-5}{}}
\newlabel{coherent-lemma-quasi-coherent-affine-cohomology-zero@cref}{{[originallemma][0][]01XB}{[1][0][]source}{}{}{}}
\makeatother
\crefname{originallemma}{original lemma}{original lemmas}
\Crefname{originallemma}{Original lemma}{Original lemmas}
\makeatletter
\newlabel{cohomology-lemma-cech-complex-complex-computes}{{0FLH}{source}{}{sourcepointer-6}{}}
\newlabel{cohomology-lemma-cech-complex-complex-computes@cref}{{[originallemma][0][]0FLH}{[1][0][]source}{}{}{}}
\makeatother
\crefname{originallemma}{original lemma}{original lemmas}
\Crefname{originallemma}{Original lemma}{Original lemmas}
\makeatletter
\newlabel{coherent-lemma-separated-case-relative-cech}{{01XL}{source}{}{sourcepointer-7}{}}
\newlabel{coherent-lemma-separated-case-relative-cech@cref}{{[originallemma][0][]01XL}{[1][0][]source}{}{}{}}
\makeatother
\crefname{originallemma}{original lemma}{original lemmas}
\Crefname{originallemma}{Original lemma}{Original lemmas}
\makeatletter
\newlabel{lemma-cohomology-base-change}{{08EU}{source}{}{sourcepointer-8}{}}
\newlabel{lemma-cohomology-base-change@cref}{{[originallemma][0][]08EU}{[1][0][]source}{}{}{}}
\makeatother
\crefname{originallemma}{original lemma}{original lemmas}
\Crefname{originallemma}{Original lemma}{Original lemmas}
\makeatletter
\newlabel{lemma-K-flat}{{0FMQ}{source}{}{sourcepointer-9}{}}
\newlabel{lemma-K-flat@cref}{{[originallemma][0][]0FMQ}{[1][0][]source}{}{}{}}
\makeatother
\crefname{originalequation}{original equation}{original equations}
\Crefname{originalequation}{Original equation}{Original equations}
\makeatletter
\newlabel{equation-kunneth-single-sheaves}{{0G49}{source}{}{sourcepointer-10}{}}
\newlabel{equation-kunneth-single-sheaves@cref}{{[originalequation][0][]0G49}{[1][0][]source}{}{}{}}
\makeatother
\crefname{originallemma}{original lemma}{original lemmas}
\Crefname{originallemma}{Original lemma}{Original lemmas}
\makeatletter
\newlabel{coherent-lemma-relative-affine-cohomology}{{089W}{source}{}{sourcepointer-11}{}}
\newlabel{coherent-lemma-relative-affine-cohomology@cref}{{[originallemma][0][]089W}{[1][0][]source}{}{}{}}
\makeatother
\crefname{originallemma}{original lemma}{original lemmas}
\Crefname{originallemma}{Original lemma}{Original lemmas}
\makeatletter
\newlabel{morphisms-lemma-affine-permanence}{{01SG}{source}{}{sourcepointer-12}{}}
\newlabel{morphisms-lemma-affine-permanence@cref}{{[originallemma][0][]01SG}{[1][0][]source}{}{}{}}
\makeatother
\crefname{originallemma}{original lemma}{original lemmas}
\Crefname{originallemma}{Original lemma}{Original lemmas}
\makeatletter
\newlabel{morphisms-lemma-pushforward-flat-affine}{{0FLM}{source}{}{sourcepointer-13}{}}
\newlabel{morphisms-lemma-pushforward-flat-affine@cref}{{[originallemma][0][]0FLM}{[1][0][]source}{}{}{}}
\makeatother
\title{The Kunneth map for a product of two flat quasi-coherent sheaves over an affine base is an isomorphism}
\author{The Stacks Project Authors}
\date{stacks-project:perfect.tex:local-snapshot-20261012; GNU FDL 1.2 or later}
\begin{document}
\maketitle
\setcounter{subsection}{1}
\noindent\textit{Editorial note.} All delivered mathematical text is the author's own, copied from the pinned Stacks Project chapter file perfect.tex; no statement or proof step was rewritten, reordered or removed. The counted span is the author environment and its own proof environment. References inside the span are handed to the assembler as navigation stubs that carry the canonical tag of the cited tag only; whether each of them is load-bearing is recorded on the stub itself.\par\medskip
\noindent\textit{Editorial note.} The last sentence of the counted proof reads, verbatim, 'This finishes the proof (except that we should check that the isomorphism is indeed given by cup product in degree 0; we omit the verification).' The omitted clause is the identification of the already constructed isomorphism with the cup product map in degree zero; the author notes it is not verified. Per the project rule the sentence is recorded here verbatim; the proof is otherwise complete and written out over more than two hundred lines.\par\medskip
\noindent\textit{Allowed standard prerequisites (editorial).}\begin{itemize}
\item The Cech complex of an affine open covering computes the cohomology of a quasi-coherent sheaf, and for a separated morphism the sheafified Cech complex represents the higher direct images. invoked at proof lines 6025-6045 and again at proof lines 6170-6175. Identifies the Cech complexes with the derived global sections and with the higher direct images of the product sheaf.
\item On a common refinement there are maps of Cech complexes compatible with pullback on cohomology, and a map of complexes from the total tensor product to the Cech complex of the tensor product compatible with the cup product. used at proof lines 5982-6018. Produces the comparison map whose bijectivity is the statement.
\item The Kunneth formula on the level of Cech complexes splits the total complex of the tensor product of the two Cech complexes as the tensor product of the cohomologies. invoked at proof lines 6060-6140. Reduces the statement to a computation with the two Cech complexes separately.
\item The Leray spectral sequence for the second projection and the base change statement for the map to the affine base compute the cohomology of the tensor product sheaf on the product. invoked at proof lines 6176-6190. Computes the cohomology of the product in terms of the cohomology of the two factors.
\end{itemize}\medskip
\ifdefined\hypertarget\hypertarget{sourcepointer-1}{}\fi
\noindent\textit{Source reference.} Original lemma 0BED: sources/stacks-chow/tags.txt, source lines 1--1. Cited in the opening sentence at proof lines 5980-5982 as reading advice; the proof of the present statement is written out in full inside the counted span.\par\medskip
\ifdefined\hypertarget\hypertarget{sourcepointer-2}{}\fi
\noindent\textit{Source reference.} Original section 01FP: sources/stacks-chow/tags.txt, source lines 1--1. Cited at proof lines 5989-5994 for the maps of Cech complexes attached to a refinement; the comparison map built from them is inside the counted span.\par\medskip
\ifdefined\hypertarget\hypertarget{sourcepointer-3}{}\fi
\noindent\textit{Source reference.} Original equation 07MB: sources/stacks-chow/tags.txt, source lines 1--1. Cited at proof lines 6000-6005 for the displayed map; the combination with the refinement maps is inside the counted span.\par\medskip
\ifdefined\hypertarget\hypertarget{sourcepointer-4}{}\fi
\noindent\textit{Source reference.} Original lemma 0FKV: sources/stacks-chow/tags.txt, source lines 1--1. Cited at proof lines 6006-6008 for compatibility with the cup product; the identification of the constructed map with the map of the statement is inside the counted span.\par\medskip
\ifdefined\hypertarget\hypertarget{sourcepointer-5}{}\fi
\noindent\textit{Source reference.} Original lemma 01XB: sources/stacks-chow/tags.txt, source lines 1--1. Cited at proof lines 6025-6030 for the identification of the Cech complexes with derived global sections; the comparison argument using it is inside the counted span.\par\medskip
\ifdefined\hypertarget\hypertarget{sourcepointer-6}{}\fi
\noindent\textit{Source reference.} Original lemma 0FLH: sources/stacks-chow/tags.txt, source lines 1--1. Cited with the previous label at proof lines 6025-6035 for the same identification; every consequent step is inside the counted span.\par\medskip
\ifdefined\hypertarget\hypertarget{sourcepointer-7}{}\fi
\noindent\textit{Source reference.} Original lemma 01XL: sources/stacks-chow/tags.txt, source lines 1--1. Cited at proof lines 6170-6175 to conclude that the sheafification of the relative Cech complex represents the higher direct images of the product sheaf; the surrounding computation is inside the counted span.\par\medskip
\ifdefined\hypertarget\hypertarget{sourcepointer-8}{}\fi
\noindent\textit{Source reference.} Original lemma 08EU: sources/stacks-perfect/pinned-original-source.tex, source lines 5269--5269. Used at proof lines 6180-6185 to compute the cohomology of the product sheaf from the affine base; the Leray step and the final conclusion are inside the counted span.\par\medskip
\ifdefined\hypertarget\hypertarget{sourcepointer-9}{}\fi
\noindent\textit{Source reference.} Original lemma 0FMQ: sources/stacks-perfect/pinned-original-source.tex, source lines 5898--5898. Used at proof lines 6135-6140 inside the Cech level computation of the higher Tor groups; the computation itself is inside the counted span.\par\medskip
\ifdefined\hypertarget\hypertarget{sourcepointer-10}{}\fi
\noindent\textit{Source reference.} Original equation 0G49: sources/stacks-perfect/pinned-original-source.tex, source lines 5949--5949. The statement refers to the displayed map by this label; the proof reconstructs that map and identifies it with the one built from Cech complexes inside the counted span.\par\medskip
\ifdefined\hypertarget\hypertarget{sourcepointer-11}{}\fi
\noindent\textit{Source reference.} Original lemma 089W: sources/stacks-chow/tags.txt, source lines 1--1. A navigation pointer used while comparing the Cech complexes on affines; the comparison itself is inside the counted span.\par\medskip
\ifdefined\hypertarget\hypertarget{sourcepointer-12}{}\fi
\noindent\textit{Source reference.} Original lemma 01SG: sources/stacks-chow/tags.txt, source lines 1--1. A navigation pointer for the affineness of the products in the induced covering; the rest of the argument is inside the counted span.\par\medskip
\ifdefined\hypertarget\hypertarget{sourcepointer-13}{}\fi
\noindent\textit{Source reference.} Original lemma 0FLM: sources/stacks-chow/tags.txt, source lines 1--1. A navigation pointer used when passing flatness of the sections through the affine projection; the use made of that flatness is inside the counted span.\par\medskip
\begin{lemma}
\label{lemma-kunneth-single-sheaf}
In the situation above the map (\ref{equation-kunneth-single-sheaves}) is an
isomorphism if $S$ is affine, $\mathcal{F}$ and $\mathcal{G}$ are $S$-flat and
quasi-coherent and $X$ and $Y$ are quasi-compact with affine diagonal.
\end{lemma}

\begin{proof}
We strongly urge the reader to read the proof of
Varieties, Lemma \ref{varieties-lemma-kunneth} first.
Choose finite affine open coverings
$\mathcal{U} : X = \bigcup_{i \in I} U_i$ and
$\mathcal{V} : Y = \bigcup_{j \in J} V_j$.
This determines an affine open covering
$\mathcal{W} : X \times_S Y = \bigcup_{(i, j) \in I \times J} U_i \times_S V_j$.
Note that $\mathcal{W}$ is a refinement of
$\text{pr}_1^{-1}\mathcal{U}$ and of $\text{pr}_2^{-1}\mathcal{V}$.
Thus by the discussion in Cohomology, Section
\ref{cohomology-section-cech-cohomology-of-complexes}
we obtain maps
$$
\check{\mathcal{C}}^\bullet(\mathcal{U}, \mathcal{F})
\to
\check{\mathcal{C}}^\bullet(\mathcal{W}, p^*\mathcal{F})
\quad\text{and}\quad
\check{\mathcal{C}}^\bullet(\mathcal{V}, \mathcal{G})
\to
\check{\mathcal{C}}^\bullet(\mathcal{W}, q^*\mathcal{G})
$$
well defined up to homotopy and compatible with pullback maps on cohomology.
In Cohomology, Equation (\ref{cohomology-equation-needs-signs})
we have constructed a map of complexes
$$
\text{Tot}(
\check{\mathcal{C}}^\bullet(\mathcal{W}, p^*\mathcal{F})
\otimes_A
\check{\mathcal{C}}^\bullet(\mathcal{W}, q^*\mathcal{G}))
\longrightarrow
\check{\mathcal{C}}^\bullet(\mathcal{W},
p^*\mathcal{F} \otimes_{\mathcal{O}_{X \times_S Y}}
q^*\mathcal{G})
$$
which is compatible with the cup product on cohomology by
Cohomology, Lemma \ref{cohomology-lemma-diagrams-commute}.
Combining the above we obtain a map of complexes
\begin{equation}
\label{equation-kunneth-on-cech}
\text{Tot}(
\check{\mathcal{C}}^\bullet(\mathcal{U}, \mathcal{F})
\otimes_A
\check{\mathcal{C}}^\bullet(\mathcal{V}, \mathcal{G}))
\to
\check{\mathcal{C}}^\bullet(\mathcal{W},
p^*\mathcal{F}
\otimes_{\mathcal{O}_{X \times_S Y}}
q^*\mathcal{G})
\end{equation}
We claim this is the map in the statement of the lemma, i.e.,
the source and target of this arrow are the same as the source
and target of (\ref{equation-kunneth-single-sheaves}). Namely, by
Cohomology of Schemes, Lemma
\ref{coherent-lemma-quasi-coherent-affine-cohomology-zero}
and
Cohomology, Lemma \ref{cohomology-lemma-cech-complex-complex-computes}
the canonical maps
$$
\check{\mathcal{C}}^\bullet(\mathcal{U}, \mathcal{F})
\to
R\Gamma(X, \mathcal{F}),
\quad
\check{\mathcal{C}}^\bullet(\mathcal{V}, \mathcal{G})
\to
R\Gamma(Y, \mathcal{G})
$$
and
$$
\check{\mathcal{C}}^\bullet(\mathcal{W},
p^*\mathcal{F} \otimes_{\mathcal{O}_{X \times_S Y}} q^*\mathcal{G})
\to
R\Gamma(X \times_S Y,
p^*\mathcal{F} \otimes_{\mathcal{O}_{X \times_S Y}}
q^*\mathcal{G})
$$
are isomorphisms. On the other hand, the complex
$\check{\mathcal{C}}^\bullet(\mathcal{U}, \mathcal{F})$
is K-flat by Lemma \ref{lemma-K-flat} and we conclude that
$\text{Tot}(
\check{\mathcal{C}}^\bullet(\mathcal{U}, \mathcal{F})
\otimes_A
\check{\mathcal{C}}^\bullet(\mathcal{V}, \mathcal{G}))$
represents the derived tensor product
$R\Gamma(X, \mathcal{F}) \otimes_A^\mathbf{L} R\Gamma(Y, \mathcal{G})$
as claimed.

\medskip\noindent
We still have to show that (\ref{equation-kunneth-on-cech})
is a quasi-isomorphism. We will do this using dimension shifting.
Set $d(\mathcal{F}) = \max \{d \mid H^d(X, \mathcal{F}) \not = 0\}$.
Assume $d(\mathcal{F}) > 0$. Set $U = \coprod\nolimits_{i \in I} U_i$.
This is an affine scheme as $I$ is finite. Denote
$j : U \to X$ the morphism which is the inclusion $U_i \to X$
on each $U_i$. Since the diagonal of $X$ is affine, the morphism
$j$ is affine, see
Morphisms, Lemma \ref{morphisms-lemma-affine-permanence}.
It follows that $\mathcal{F}' = j_*j^*\mathcal{F}$ is $S$-flat, see
Morphisms, Lemma \ref{morphisms-lemma-pushforward-flat-affine}.
It also follows that $d(\mathcal{F}') = 0$ by combining
Cohomology of Schemes, Lemmas
\ref{coherent-lemma-relative-affine-cohomology} and
\ref{coherent-lemma-quasi-coherent-affine-cohomology-zero}.
For all $x \in X$ we have $\mathcal{F}_x  \to \mathcal{F}'_x$
is the inclusion of a direct summand: if $x \in U_i$,
then $\mathcal{F}' \to (U_i \to X)_*\mathcal{F}|_{U_i}$
gives a splitting. We conclude that
$\mathcal{F} \to \mathcal{F}'$ is injective and
$\mathcal{F}'' = \mathcal{F}'/\mathcal{F}$
is $S$-flat as well. The short exact sequence
$0 \to \mathcal{F} \to \mathcal{F}' \to \mathcal{F}'' \to 0$
of flat quasi-coherent $\mathcal{O}_X$-modules
produces a short exact sequence of complexes
$$
0 \to
\text{Tot}(
\check{\mathcal{C}}^\bullet(\mathcal{U}, \mathcal{F})
\otimes_A
\check{\mathcal{C}}^\bullet(\mathcal{V}, \mathcal{G})) \to
\text{Tot}(
\check{\mathcal{C}}^\bullet(\mathcal{U}, \mathcal{F}')
\otimes_A
\check{\mathcal{C}}^\bullet(\mathcal{V}, \mathcal{G})) \to
\text{Tot}(
\check{\mathcal{C}}^\bullet(\mathcal{U}, \mathcal{F}'')
\otimes_A
\check{\mathcal{C}}^\bullet(\mathcal{V}, \mathcal{G})) \to 0
$$
and a short exact sequence of complexes
$$
0 \to
\check{\mathcal{C}}^\bullet(\mathcal{W},
p^*\mathcal{F}
\otimes_{\mathcal{O}_{X \times_S Y}}
q^*\mathcal{G}) \to
\check{\mathcal{C}}^\bullet(\mathcal{W},
p^*\mathcal{F}'
\otimes_{\mathcal{O}_{X \times_S Y}}
q^*\mathcal{G}) \to
\check{\mathcal{C}}^\bullet(\mathcal{W},
p^*\mathcal{F}''
\otimes_{\mathcal{O}_{X \times_S Y}}
q^*\mathcal{G}) \to 0
$$
Moreover, the maps (\ref{equation-kunneth-on-cech}) between these are
compatible with these short exact sequences. Hence it suffices to prove
(\ref{equation-kunneth-on-cech})
is an isomorphism for $\mathcal{F}'$ and $\mathcal{F}''$. Finally,
we have $d(\mathcal{F}'') < d(\mathcal{F})$.
In this way we reduce to the case $d(\mathcal{F}) = 0$.

\medskip\noindent
Arguing in the same fashion for $\mathcal{G}$ we find that we
may assume that both $\mathcal{F}$ and $\mathcal{G}$
have nonzero cohomology only in degree $0$.
Observe that this means that $\Gamma(X, \mathcal{F})$
is quasi-isomorphic to the $K$-flat complex
$\check{\mathcal{C}}^\bullet(\mathcal{U}, \mathcal{F})$
of $A$-modules sitting in degrees $\geq 0$.
It follows that $\Gamma(X, \mathcal{F})$ is a flat $A$-module
(because we can compute higher Tor's against this module
by tensoring with the Cech complex).
Let $V \subset Y$ be an affine open. Consider the affine open covering
$\mathcal{U}_V : X \times_S V = \bigcup_{i \in I} U_i \times_S V$.
It is immediate that
$$
\check{\mathcal{C}}^\bullet(\mathcal{U}, \mathcal{F})
\otimes_A \mathcal{G}(V) =
\check{\mathcal{C}}^\bullet(\mathcal{U}_V,
p^*\mathcal{F} \otimes_{\mathcal{O}_{X \times Y}}
q^*\mathcal{G})
$$
(equality of complexes). By the flatness of $\mathcal{G}(V)$
over $A$ we see that
$\Gamma(X, \mathcal{F}) \otimes_A \mathcal{G}(V) \to
\check{\mathcal{C}}^\bullet(\mathcal{U}, \mathcal{F})
\otimes_A \mathcal{G}(V)$ is a quasi-isomorphism.
Since the sheafification of
$V \mapsto \check{\mathcal{C}}^\bullet(\mathcal{U}_V,
p^*\mathcal{F} \otimes_{\mathcal{O}_{X \times Y}}
q^*\mathcal{G})$ represents
$Rq_*(p^*\mathcal{F} \otimes_{\mathcal{O}_{X \times Y}} q^*\mathcal{G})$
by Cohomology of Schemes, Lemma
\ref{coherent-lemma-separated-case-relative-cech}
we conclude that
$$
Rq_*(p^*\mathcal{F} \otimes_{\mathcal{O}_{X \times Y}} q^*\mathcal{G})
\cong
\Gamma(X, \mathcal{F}) \otimes_A \mathcal{G}
$$
on $Y$ where the notation on the right hand side indicates the module
$$
b^*\widetilde{\Gamma(X, \mathcal{F})} \otimes_{\mathcal{O}_Y} \mathcal{G}
$$
Using the Leray spectral sequence for $q$ we find
$$
H^n(X \times_S Y, p^*\mathcal{F} \otimes_{\mathcal{O}_{X \times Y}}
q^*\mathcal{G}) =
H^n(Y,
b^*\widetilde{\Gamma(X, \mathcal{F})} \otimes_{\mathcal{O}_Y} \mathcal{G})
$$
Using Lemma \ref{lemma-cohomology-base-change} for the morphism
$b : Y \to S = \Spec(A)$ and using that $\Gamma(X, \mathcal{F})$
is $A$-flat we conclude that
$H^n(X \times_S Y, p^*\mathcal{F} \otimes_{\mathcal{O}_{X \times Y}}
q^*\mathcal{G})$ is zero for $n > 0$ and isomorphic to
$H^0(X, \mathcal{F}) \otimes_A H^0(Y, \mathcal{G})$ for $n = 0$.
Of course, here we also use that $\mathcal{G}$ only has
cohomology in degree $0$.
This finishes the proof (except that we should check that the
isomorphism is indeed given by cup product in degree $0$; we omit
the verification).
\end{proof}

\end{document}
